\bigskip
\bigskip
\subsubsection*{Խնդիրներ և հարցեր թեմա 14-ի վերաբերյալ}

\begin{enumerate}[label=\thesection.\arabic*.]

% 15.1
\item Ճի՞շտ է պնդումը․ կապակցված տարածության ցանկացած ենթաբազ\-մու\-թյուն կապակցված է։

% 15.2
\item Ապացուցեք․ կապակցված $X$ տարածության ցանկացած $\faktor{X}{R}$ ֆակտոր-տարա\-ծություն կապակցված է։

% 15.3
\item Ապացուցեք․ $(\R, \text{սովոր.})$ տարածության ռացիոնալ կետերի ամեն մի ենթա\-բազ\-մություն կապակցված չէ։

% 15.4
\item Ունենք $f: X \to Y$ անընդհատ արտապատկերում։ Ճի՞շտ է պնդումը․ $Y$-ում կա\-պա\-կցված ամեն մի $B$ ենթաբազմության $f^{-1}(B)$ նախակերպարը կապա\-կց\-ված է $X$-ում։

% 15.5
\item Ապացուցեք․ $\R^n$ էվկլիդեսյան տարածության ամեն մի ուռուցիկ ենթա\-բազ\-մություն կապակցված ենթատարածություն է:
\begin{hint}
Սևեռելով $X \subset \R^n$ ուռուցիկ ենթաբազմության որևէ $x_0$ կետ՝ դի\-տար\-կեք $x_0$ սկզբնակետով բոլոր հատվածները, որոնք ընկած են $X$-ում:
\end{hint}

% 15.6
\item $\R^3$ էվկլիդեսյան տարածության հետևյալ երեք՝
\[
\big\{(x,y,z); \; x^2+y^2-z^2=\delta\big\}, 
\qquad \text{որտեղ } \delta=-1, 0, 1
\] 
մակերևույթներից որո՞նք են կապակցված (պատասխանը հիմնավորե՛ք):

% 15.7
\item Ճի՞շտ է արդյոք պնդումը․ $(\R, \mapsto)$ տարածության հետևյալ ենթա\-բազ\-մու\-թյուն\-ներից՝ ա) $\{0,1\}$, բ) $[0,1)$, գ) $[0,1]$, դ) $(0,0)$, ե) $(0,1]$, ոչ մեկը կապակցված չէ:

% 15.8
\item Ապացուցեք․ $\R^2$ էվկլիդեսյան հարթության այն բոլոր կետերի $A$ ենթա\-բազ\-մու\-թյունը, որոնց կոորդինատներից գոնե մեկը ռացիոնալ թիվ է, կապակցված ենթա\-բազմություն է:
\begin{hint}
Հիմնավորեք, որ $A$-ի ցանկացած երկու կետեր կարող են միացվել իրար երկու կամ երեք հատվածներից կազմված բեկյալով, որն ամբողջությամբ ընկած է $A$-ի մեջ:
\end{hint}

% 15.9
\item Ճի՞շտ է արդյոք, որ $\R^2$ հարթության այն բոլոր կետերի $B$ ենթաբազմությունը, որոնց կոորդինատներից միայն մեկն է ռացիոնալ թիվ, կապակցված ենթա\-բազ\-մություն է:
\begin{hint}
Ըստ պայմանի՝ $B$-ն չունի ընդհանուր կետ $y=x$ ուղղի հետ։ Ուստի $B$-ն ներկայացվում է որպես երկու ենթաբազմությունների միավորում, որոնցից մեկը ընկած է այդ ուղղից վերև, իսկ մյուսը՝ ներքև:
\end{hint}

% 15.10
\item Ապացուցեք․ $\R$ թվային ուղղի որևէ $[a,b]$ ենթատարածություն հոմեոմորֆ չէ $\R^2$ հարթության որևէ $[c,d] \times [e,f]$ ենթատարածության:

% 15.11
\item Ապացուցեք․ $X$ տոպոլոգիական տարածությունը կապակցված է այն և միայն այն դեպքում, երբ ցանկացած $f: X \to Y$ անընդհատ արտապատկերում, որ\-տեղ $Y$-ը մեկից ավելի կետեր պարունակող դիսկրետ տարածություն է, հաս\-տա\-տուն արտապատկերում է:

% 15.12
\item Ապացուցեք․ $(X, \tau)$ տոպոլոգիական տարածությունը կապակցված չէ այն և մի\-այն այն դեպքում, երբ գոյություն ունի $f: (X, \tau) \to (Y, \text{դիսկր.})$ անընդհատ սյուրյեկտիվ արտապատկերում, որտեղ $Y$-ը երկկետանոց բազմություն է:

% 15.13
\item Ապացուցեք․ եթե $A$-ն $X$ տարածության կապակցված ենթատարածություն է և $A \subset Y \subset \overline{A}$, ապա $Y$-ը $X$-ի կապակցված ենթատարածություն է:

\end{enumerate}
